\section{Formal Reinterpretation of the Concept-Injection Experiments}

This chapter uses the theorems of Chapter 3 to systematically reinterpret the empirical
results of arXiv:2601.01828\cite{lindsey-2026-introspection}. There are two aims: first, to
test the falsifiability of the theorems (if the paper's data contradicted the theorems'
predictions, the theorems would be refuted); second, to show that the ``weak functional
self-report'' observed by the paper is \textbf{not} evidence of in-paradigm introspective
capacity but a necessary corollary of the theorems. All cited data come from explicit
statements in the paper's main text and appendix.

\subsection{The Injected-Thought-Detection Experiment: Why 20\% Success Is a Capability Ceiling, Not a Structural Guarantee}

The paper reports: Claude Opus 4.1 and Opus 4 detect injected thoughts at rates of about 20\%
under the optimal injection layer and strength; production models produce zero false positives
over 100 control trials; random vectors elicit ``noticing'' only 9 times out of 100
trials\cite{lindsey-2026-introspection}.

Reading via Theorem 1b: \textbf{the training objective provides no guarantee of introspective
accuracy}, so 20\% is not evidence of ``capability masked by noise'' but should be understood as
a by-product of ``inductive bias happening to cover part of the query space''. Concretely:
detecting an injected concept requires the model to implement an ``anomaly-detection function''
(one of the mechanism conjectures the paper lists in Section 10.3.1
\cite{lindsey-2026-introspection}). That function can \textbf{possibly} be learned because it
is useful as a \textbf{by-product of text prediction} (predicting anomalous text and detecting
implausible content are common corpus tasks), not because the training objective rewards
``faithful reports about introspective states''. Theorem 1a further shows that, in the
\textbf{unintervened} native regime, the report--state mutual information is zero; the 20\%
success depends entirely on the experimenter's external injection (artificially altering the
state and letting the model ``read'' the altered state).

The paper's 20\% is therefore an incidental capability ceiling, attained under \textbf{external
intervention} and on a \textbf{specific query distribution}, with \textbf{no guarantee} of any
kind; it is far removed from the assertion of ``in-paradigm introspective capacity''. The paper's
own phrasing (``highly unreliable and context-dependent'', ``failures of
introspection remain the norm''\cite{lindsey-2026-introspection}) agrees completely with the
predictions of
Theorem 1.

\subsection{Distinguishing ``Thoughts'' from Text: The Public Computability of Consistency-Check Circuits}

The paper's second experiment shows that, under some conditions, models can simultaneously
report the injected concept and transcribe the input text, with all models performing
significantly above chance\cite{lindsey-2026-introspection}. The paper attributes this to
different attention heads retrieving different kinds of information (Section 10.3.2).

Reading via Theorem 5: \textbf{this capacity is publicly computable}. The process by which the
model ``distinguishes thoughts from text'' is a deterministic function of $(\theta, c, u)$
($u$ the injected vector, known to the experimenter), so any third party holding these
quantities can reproduce the same discrimination. This accords with the paper's own caveat
that ``the mechanism may be shallow and narrow''\cite{lindsey-2026-introspection}. More
importantly: even if the model \textbf{behaviorally} distinguishes ``thoughts'' from ``text'',
the basis of its discrimination (the internal attention pattern) lies entirely within the
read range of an external probe; there is no private channel in the information-theoretic
sense, and the privileged-access criterion\cite{song-privileged-2025} fails (Theorem 5b).

\subsection{The Prefill-Detection Experiment: The Statistical Nature of Causal Attribution}

The paper's most dramatic result is the ``bread'' experiment\cite{lindsey-2026-introspection}:
when the model's output is forcibly changed to ``bread'', it denies that the output was
intentional; but when the experimenter injects the ``bread'' concept vector into the activations
before the prefill, the model instead admits that the output was intentional and
\textbf{fabricates after the fact} a rationalization.

Joint reading via Theorems 4 and 5: the model's ``intentionality judgment'' is a statistical
adjudication by a \textbf{consistency-check circuit} (the paper's Section 10.3.3 cites the
concordance-heads mechanism\cite{kamath-2025-concordance}) over $(\theta, c, u)$, comparing
``the intentionality representation in prior activations'' with ``the likelihood of the actual
output''. Injecting the ``bread'' vector changes the statistical input and therefore flips the
judgment. This corresponds precisely to the key distinction emphasized repeatedly
above: ``\textbf{statistical causal association, not phenomenological subjective
experience}''. The model does not ``experience'' intentionality; it merely runs a statistical
classifier (one that injection can fool). The fabricated post-hoc rationalization is exactly
what Theorem 4 predicts: the model cannot verify the grounding of its own statements (no
verifier node) and can only invoke ``explanation text patterns'' learned from the corpus to
complete the narrative.

\subsection{Zero Net Benefit for Base Models: Post-Training Is Only an Optimization of Text Patterns}

The paper reports that base (pretrained) models generally have high false-positive rates on
the injection-detection task and net benefit (true-positive rate minus false-positive rate)
at or below zero; post-training strategies significantly affect introspective-task
performance, with production models ``reluctant to participate in introspective exercises''
while refusal-avoiding (``helpful-only'') variants perform better\cite{lindsey-2026-introspection}.

Reading via Theorem 2: the pretraining objective is pure text likelihood, meaning statistical
learning of ``introspective-style text''; base models treat both ``detecting the injection'' and
``saying something'' as text patterns, hence high false positives and poor grounding (zero net
benefit). The reward signals introduced at the post-training stage (RLHF-type) are still, in
essence, shaping of \textbf{textual behavior} (answering like an ``introspective assistant''),
not supervision of ``faithful state reporting''; they merely redistribute probability mass
over text patterns. The paper's observation that ``helpful-only variants perform better''
shows precisely that what is optimized is never \textbf{grounding} but \textbf{textual
behavior that looks like introspection}. This is exactly the prediction of Theorem 2b:
optimizing a textual objective can systematically deviate from faithful state reporting.

\subsection{Layer Sensitivity: Multiple Narrow Circuits, Not a Unified Introspective Mechanism}

The paper finds that different introspective behaviors have different sensitivities to the
injection layer (the optimal layer for detecting injected thoughts lies at about two-thirds
depth, while that for prefill detection is shallower)\cite{lindsey-2026-introspection}, and
infers that ``introspection is not supported by a single mechanism''.

Reading via Theorem 1b: since the training objective does not reward a unified introspective
mechanism, any observed introspective-like behavior must be a by-product of \textbf{several
independent, narrow circuits that serve text prediction}, whose layer distributions are
naturally different. This agrees completely with the paper's Section 10.3 ``simplest mechanism''
conjectures (anomaly-detection circuits, retrieval heads, consistency-check circuits,
salience-marking circuits\cite{lindsey-2026-introspection}), and provides a theoretical basis
for the \textbf{necessity} of these conjectures: the paradigm does not reward a unified
mechanism, so what is observed must be an assemblage of circuits.

\subsection{The Paper's Self-Reported Limitations versus the Theorems}

\begin{table}[htbp]
\centering
\caption{Correspondence between the paper's self-reported limitations and the theorems
\label{tab:paper-thm}}
\begin{tabular}{p{7.2cm}p{6.2cm}}
\toprule
Self-reported limitation (quoted) & Corresponding theorem \\
\midrule
``The abilities we observe are highly unreliable'' & Theorem 1b: no PAC guarantee of any kind \\
``Failures of introspection remain the norm'' & Theorem 1a: zero mutual information in the native regime \\
``the mechanisms underlying our results could still be rather shallow and narrowly specialized'' & Theorem 1b + Section 4.5: narrow-circuit by-products \\
``Our concept injection protocol places models in an unnatural setting unlike those they face in training or deployment'' & Theorem 1a: capacity depends on external intervention \\
``Demonstrating metacognitive representations is difficult to do directly, and we do not do so in this work'' & Theorem 4: the generation objective does not distinguish with/without metarepresentation \\
``we do not seek to address the question of whether AI systems possess human-like self-awareness or subjective experience'' & Theorem 5: no privileged channel (first-person access unavailable) \\
\bottomrule
\end{tabular}
\end{table}

\begin{remark}[Falsifiability statement]
The consistency between the theorems of this report and the paper's data is not circular: the
derivation of the theorems depends only on the postulates $\Pi_1$--$\Pi_5$ and standard
mathematical results (information theory, NFL, Goodhart), not on any of the paper's
experimental data; the paper's data serve here only as \textbf{tests}. If a future system
satisfying all postulates stably achieved $I(R;H\mid\theta,c) > 0$ in the native regime
(e.g., a reliable description of its own activations without any intervention), Theorem 1a
would be refuted. We consider this impossible within the
postulates, but we would welcome such experimental evidence.
\end{remark}

\clearpage
